% 直近の変更履歴（課題1・課題2 用）
% management_plan.tex の \chapter{変更履歴} から \input される．
%
% ライフサイクル:
%   - 変更があったら，このファイルに行を追記する（書式は下記）
%   - 提出のタイミングで，下記の履歴行を revision_history/sections/changelog.tex
%     （累積版・課題3用）の末尾にコピーしてから，このファイルの履歴行を全て削除する
%
% 書式:
%   YYYY-MM-DD & 学籍番号 : 氏名 & 対象（章・節） & 変更内容 \\
%
% 前提パッケージ（呼び出し側でロード済みであること）:
%   \usepackage{longtable}
%   \usepackage{booktabs}
%   \usepackage{array}

\begin{longtable}{p{2.2cm}p{3.5cm}p{4cm}p{5.5cm}}
\toprule
日付 & 担当 & 対象（章・節） & 変更内容 \\
\midrule
\endfirsthead

\toprule
日付 & 担当 & 対象（章・節） & 変更内容 \\
\midrule
\endhead

% --- 履歴行をここに追記していく（提出時にクリア） ---
2026-07-01 & 24G1056 : 慶田 優 & 第6章 テスト実施記録 & INT-5（カラーセンサ＋フォトトランジスタ統合テスト）の結果を追記（不合格：環境光によるS13683認識不安定） \\
2026-07-01 & 24G1056 : 慶田 優 & 第5章 TPM測定推移 & TPM-2（拍検出成功率）を3時点分追記（6/17単体≈100\%○，6/24結合≈90\%×，7/1測定不能×）．TPM-4（EMA収束拍数：約1拍○）・TPM-5（最大追従BPM：110bpm○）を追記 \\
2026-07-01 & 24G1056 : 慶田 優 & 第4章 TRL評価 & 拍タイミング検出の根拠にS13683環境光問題を追記．テンポ推定の根拠を「3拍以内」から「約1拍以内（TPM-4実測値）」に更新．テンポ追従演奏（TRL1→3），カラーLED制御（輝度エンコードLEDから改称・根拠更新），演奏状態遷移管理（TRL2→3）を更新 \\

\bottomrule
\endfoot
\end{longtable}
